\documentclass[popl.tex]{subfiles}
\begin{document}

\clearpage\section{Example Repairs}\label{sec:exaple_repairs}

Below, we provide a few representative examples of broken code snippets and the corresponding human repairs that were successfully ranked first by our method. On the left is a complete snippet fed to the model, and on the right, the corresponding human repair that was correctly predicted.

\begin{figure}[H]
\begin{tabular}{|m{6.6cm}|m{6.6cm}|}
\hline \rule{0pt}{2.5ex}\textbf{Original broken code}\rule[-1ex]{0pt}{2ex} &  \rule{0pt}{2.5ex}\textbf{First predicted repair}\rule[-1ex]{0pt}{2ex} \\\hline
\begin{smallpy}

(*@\hlorange{form}@*) sympy import *
x = Symbol('x', real=True)
x, re(x), im(x)

\end{smallpy} & \begin{smallpy}

(*@\hlorange{\textbf{from}}@*) sympy import *
x = Symbol('x', real=True)
x, re(x), im(x)

\end{smallpy} \\\hline
\begin{smallpy}

result = (*@\hlorange{yeald}@*) From(item.create())
raise Return(result)

\end{smallpy} & \begin{smallpy}

result = (*@\hlorange{\textbf{yield}}@*) From(item.create())
raise Return(result)

\end{smallpy} \\\hline
\begin{smallpy}

df.apply(lambda row: list(set(row['ids'(*@\hlorange{)}@*))))

\end{smallpy} & \begin{smallpy}

df.apply(lambda row: list(set(row['ids'(*@\hlorange{]}@*))))

\end{smallpy} \\\hline
%        \begin{lstlisting}[basicstyle=\ttfamily\lst@ifdisplaystyle\footnotesize\fi, language=python]
%
%  import numpy (*@\hlorange{ad}@*) np
%  A_concate = np.array([a_0, a_1, a_2,..., a_n])
%
%        \end{lstlisting} & \begin{lstlisting}[basicstyle=\ttfamily\lst@ifdisplaystyle\footnotesize\fi, language=python]
%
%  import numpy (*@\hlorange{\textbf{as}}@*) np
%  A_concate = np.array([a_0, a_1, a_2,..., a_n])
%
%        \end{lstlisting} \\\hline
\begin{smallpy}

sum(len(v) for v items.values())(*@\hlred{)}@*)

\end{smallpy} & \begin{smallpy}

sum(len(v) for v (*@\hlgreen{\textbf{in}}@*) items.values())

\end{smallpy} \\\hline
\begin{smallpy}

def average(values):
  if values == (1,2,3):
    return (1+2+3)/3
  (*@\hlorange{else}@*) (*@\hlred{if}@*) values == (-3,2,8,-1):
    return (-3+2+8-1)/4

\end{smallpy} & \begin{smallpy}

def average(values):
  if values == (1,2,3):
    return (1+2+3)/3
  (*@\hlorange{elif}@*) values == (-3,2,8,-1):
    return (-3+2+8-1)/4

\end{smallpy} \\\hline
\begin{smallpy}

dict = {
  "Jan": 1
  "January": 1
  "Feb": 2 # and so on
}

\end{smallpy} & \begin{smallpy}

dict = {
  "Jan": 1(*@\hlgreen{,}@*)
  "January": 1(*@\hlgreen{,}@*)
  "Feb": 2 # and so on
}

\end{smallpy} \\\hline
\begin{smallpy}

class MixIn(object)
  def m():
    pass

class classA(MixIn):

class classB(MixIn):

\end{smallpy} & \begin{smallpy}

class MixIn(object)(*@\hlgreen{:}@*)
  def m():
    pass

class classA(MixIn): (*@\hlgreen{\textbf{pass}}@*)

class classB(MixIn): (*@\hlgreen{\textbf{pass}}@*)

\end{smallpy} \\\hline
\end{tabular}
\end{figure}
\ifSubfilesClassLoaded{\par\clearpage\bibliography{../bib/acmart}}{}
\end{document}
